% Ce chapitre repose sur :
%   DOC: s12, s13, s14, s20, s26, s27
%   DOC: s30
%   OBS: REL-RDV.C02, REL-RDV.C07, REL-RDV.C08, REL-RDV.R08, REL-RDV.R11
%   HYP: (aucun)
%   CHI: absences-avec-agent-annulation, absences-total, activites-distinctes, annulations-avec-agent, annulations-total, artefact-nombre
%   CHI: artefact-part-arrondi, cible-delai-niveau-1, cible-delai-niveau-2, cible-delai-niveau-3, cible-delai-niveau-4, cible-delai-niveau-5
%   CHI: code-activite-cardiologie, code-activite-dermatologie, code-activite-gyneco-obstetrique, code-activite-medecine-generale, code-activite-medecine-interne, code-activite-ophtalmologie
%   CHI: code-activite-oto-rhino-laryngologie, code-activite-pediatrie, collision-groupes-nom, collision-groupes-piece, complement-nombre, creances-part-au-dela-un-an-arrondi
%   CHI: delai-mediane-complement, delai-mediane-ensemble, delai-rapport-extremes, doublons-part-mesuree-arrondi, encaissements-montant, facturation-factures
%   CHI: facturation-montant-total, facturation-part-organisme-arrondi, facturation-part-patient-arrondi, feries-jours, feries-programme-passages, feries-urgences-hors-moyenne-arrondi
%   CHI: feries-urgences-moyenne-arrondi, flux-dimanche-programme, flux-heure-pointe, flux-heures-vides-programme, flux-jours-ouverts, flux-lundi-moyenne-arrondi
%   CHI: flux-passages-heure-pointe, flux-passages-total, flux-programme-moyenne-arrondi, flux-programme-total, flux-samedi-moyenne-arrondi, flux-urgences-moyenne
%   CHI: grappes-multiples, grappes-taille-maximale, hebdo-rapport-mesure-arrondi, honores-avec-fiche, honores-sans-fiche, honores-sans-fiche-en-quarantaine
%   CHI: honores-total, identifiants-distincts, injecte-correlation-parametres-arrondi, injecte-part-jour-meme, injecte-taux-annulation, jour-meme-marge-nulle
%   CHI: jour-meme-part-marge-nulle-arrondi, niveaux-tri, part-jour-meme-arrondi, passages-urgences-produits, periode-debut, periode-fin
%   CHI: periode-jours, periode-mois, personnes-estimees, qualite-colonne-vide, qualite-colonnes, qualite-colonnes-incompletes
%   CHI: qualite-lignes-examinees, qualite-tables, qualite-tables-sans-quarantaine, qualite-taux-quarantaine-maximal-arrondi, quarantaine-lignes-total-analyse, ramadan-baisse-arrondi
%   CHI: ramadan-jours, ramadan-jours-ouverts, rdv-en-instance, rdv-etats, rdv-etats-signification, recouvrement-creances
%   CHI: recouvrement-montant-du, recouvrement-montant-restant, recouvrement-relances, recouvrement-taux-aboutissement-arrondi, recouvrement-taux-global-arrondi, recouvrement-taux-maximal-arrondi
%   CHI: recouvrement-taux-minimal-arrondi, relation-brute-complement-arrondi, relation-brute-ensemble-arrondi, relation-effectif-complement, relation-effectif-ensemble, relation-inter-complement-arrondi
%   CHI: relation-inter-ensemble-arrondi, relation-intra-complement-arrondi, relation-intra-ensemble-arrondi, source-part-non-appropriees-arrondi, source-part-non-urgentes, source-rendez-vous-lignes
%   CHI: taux-absenteisme-global-arrondi, taux-annulation-global-arrondi, taux-encaissement-arrondi, urgences-heure-creux, urgences-heure-pointe, urgences-part-consultation-ordinaire-arrondi
%   CHI: urgences-part-hospitalisation-arrondi, urgences-part-mode-principal-arrondi, urgences-part-retour-domicile-arrondi, urgences-part-transfert-arrondi, urgences-rapport-pointe-creux-arrondi
%   SER: absenteisme-par-activite, anciennete-creances, completude-champs, delais-par-activite, delais-repartition, flux-horaire
%   SER: flux-mensuel, urgences-par-niveau
%   SECTIONS: 6

\chapter{Analyse de l'activité}\label{chap:activite}

La chaîne calcule des grandeurs sur les flux, les délais, l'absentéisme, les urgences, la
facturation et la qualité du fichier patient. La question traitée est ce que ces grandeurs montrent,
et ce qui, en elles, vient des paramètres du générateur. Les chiffres proviennent du jeu de données
généré : ils illustrent ce que la chaîne sait mesurer, et non l'activité réelle du service.
L'annexe~\ref{annexe:correspondance} indique, pour chaque conclusion, la relation injectée qu'elle
reproduit ou l'absence de relation.

\section{Les flux de patients}

Les valeurs de ce chapitre sont arrondies ; leurs valeurs complètes figurent au
tableau~\ref{tab:precision-complete}.

Les deux flux de patients diffèrent par leur rythme annuel comme par leur profil horaire. Le fait de
passage réunit \chiffre{flux-passages-total} lignes du \chiffre{periode-debut} au
\chiffre{periode-fin}, soit \chiffre{periode-jours} jours. L'activité programmée, consultation
externe et hospitalisation, en compte \chiffre{flux-programme-total}, et les urgences le reste. Ces
volumes dépendent de la calibration des passages aux urgences, dont la
section~\ref{sec:sensibilite} mesure l'effet. Les taux, les parts et les délais n'en dépendent pas.
Les deux flux se séparent par le code de type de passage, car aucune source ne documente son
libellé (voir figure~\ref{fig:passages-mois}).

\begin{figure}[htbp]
\centering
\begin{tikzpicture}
\begin{axis}[
  width=12cm, height=4.6cm,
  xlabel={année}, ylabel={passages par mois},
  ymin=0, xtick={2024,2024.5,2025,2025.5,2026,2026.5},
  xticklabels={2024, , 2025, , 2026, },
  grid=major, grid style={gray!25},
  every axis plot/.append style={thick},
  legend style={at={(0.5,-0.42)}, anchor=north, legend columns=2, font=\footnotesize,
    % SÉPARATION ENTRE LES DEUX COLONNES DE LA LÉGENDE. Sans elle, la clé de la seconde entrée
    % touchait le libellé de la première, et les deux se lisaient d'un seul tenant. Le défaut se
    % voyait sur les quatre graphiques à légende en deux colonnes, et sur eux seuls.
    /tikz/every even column/.append style={column sep=1.2em},
    draw=gray!40},
]
\addplot[color=black, mark=none] table[x=annee, y=urgences, col sep=comma] {\serie{flux-mensuel}};
\addlegendentry{urgences}
\addplot[color=black, dashed, mark=none] table[x=annee, y=programme, col sep=comma]
  {\serie{flux-mensuel}};
\addlegendentry{activité programmée}
\end{axis}
\end{tikzpicture}
\caption{Passages par mois, sur les \chiffre{periode-mois} mois de la période}\label{fig:passages-mois}
\end{figure}

\conclusion{C-01}\textbf{Le flux des urgences reste plat sur toute la période, quand l'activité
programmée varie de mois en mois.} Les urgences reçoivent \chiffre{flux-urgences-moyenne} passages
par jour, tous les jours. L'activité programmée en reçoit \chiffre{flux-programme-moyenne-arrondi}
par jour ouvert. Cette moyenne se rapporte aux \chiffre{flux-jours-ouverts} jours qui portent au
moins un passage programmé. Rapportée aux \chiffre{periode-jours} jours de la période, elle
baisserait d'un cinquième sans que l'activité change.

Le profil horaire des deux flux diffère aussi (voir figure~\ref{fig:passages-heure}).

\begin{figure}[htbp]
\centering
\begin{tikzpicture}
\begin{axis}[
  width=12cm, height=4.4cm, ybar, bar width=3.4pt,
  xlabel={heure d'entrée}, ylabel={passages},
  xmin=-0.8, xmax=23.8, ymin=0, xtick={0,3,6,9,12,15,18,21},
  grid=major, grid style={gray!25},
  legend style={at={(0.5,-0.45)}, anchor=north, legend columns=2, font=\footnotesize,
    % SÉPARATION ENTRE LES DEUX COLONNES DE LA LÉGENDE. Sans elle, la clé de la seconde entrée
    % touchait le libellé de la première, et les deux se lisaient d'un seul tenant. Le défaut se
    % voyait sur les quatre graphiques à légende en deux colonnes, et sur eux seuls.
    /tikz/every even column/.append style={column sep=1.2em},
    draw=gray!40},
]
\addplot[fill=gray!30, draw=gray!60] table[x=heure, y=urgences, col sep=comma]
  {\serie{flux-horaire}};
\addlegendentry{urgences}
\addplot[fill=black, draw=black] table[x=heure, y=programme, col sep=comma] {\serie{flux-horaire}};
\addlegendentry{activité programmée}
\end{axis}
\end{tikzpicture}
\caption{Passages par heure d'entrée, cumulés sur la période entière}\label{fig:passages-heure}
\end{figure}

\conclusion{C-02}\textbf{L'activité programmée tient dans les heures d'ouverture.} Elle culmine à
\chiffre{flux-heure-pointe} heures avec \chiffre{flux-passages-heure-pointe} passages, et
\chiffre{flux-heures-vides-programme} heures de la journée n'en portent aucun. Les urgences reçoivent
des patients à chaque heure. La même requête trouve des heures vides sur le flux programmé, ce qui
rend fiable leur absence aux urgences.

\conclusion{C-03}\textbf{La charge programmée baisse du premier au sixième jour de la semaine et
tombe à zéro le septième.} Elle passe de \chiffre{flux-lundi-moyenne-arrondi} passages par jour le
premier jour à \chiffre{flux-samedi-moyenne-arrondi} le sixième, soit un rapport de
\chiffre{hebdo-rapport-mesure-arrondi}. Le septième jour porte \chiffre{flux-dimanche-programme}
passage programmé.

\conclusion{C-04}\textbf{Le Ramadan abaisse l'activité programmée sans l'interrompre.} La période
compte \chiffre{ramadan-jours} jours de Ramadan, dont \chiffre{ramadan-jours-ouverts} portent une
activité programmée. Sur ces jours, la moyenne par jour ouvert baisse de
\chiffre{ramadan-baisse-arrondi}~\%. Les dates du Ramadan viennent d'une source \cite{s26}, et
l'ampleur de la baisse vient d'un paramètre du générateur.

\conclusion{C-05}\textbf{Un jour férié ferme l'activité programmée et laisse les urgences
inchangées.} Les \chiffre{feries-jours} jours fériés de la période portent
\chiffre{feries-programme-passages} passage programmé. Les urgences y reçoivent
\chiffre{feries-urgences-moyenne-arrondi} passages par jour, contre
\chiffre{feries-urgences-hors-moyenne-arrondi} les autres jours.

\section{Les délais d'obtention de rendez-vous}

Le délai d'obtention dépend en grande partie de sa construction, mais son centre reste utilisable.
Il est l'écart entre la date de création du rendez-vous\releve{REL-RDV.C02} et la date du
rendez-vous\releve{REL-RDV.R08}, deux champs relevés au poste.

Le générateur tire un délai dans une loi log-normale, puis le plafonne à la marge disponible. Cette
marge est le nombre de jours entre la création de la fiche patient et le rendez-vous. Quand la
marge est plus courte que le délai tiré, le délai obtenu vaut exactement la marge. Sur les
\chiffre{honores-total} rendez-vous honorés, \chiffre{honores-avec-fiche} ont une fiche patient en
vigueur. Les \chiffre{honores-sans-fiche} autres n'en ont pas, car leurs
\chiffre{honores-sans-fiche-en-quarantaine} fiches sont en quarantaine.

Parmi les \chiffre{honores-avec-fiche} rendez-vous rangés, \chiffre{artefact-nombre} portent un délai
égal à la marge, soit \chiffre{artefact-part-arrondi}~\%. Les \chiffre{complement-nombre} autres
forment le complément. Les relations qui mettent en jeu le délai sont donc mesurées deux fois, sur
l'ensemble et sur ce complément, à commencer par la répartition du délai (voir
figure~\ref{fig:delai-repartition}).

\begin{figure}[htbp]
\centering
\begin{tikzpicture}
\begin{axis}[
  width=11cm, height=4.6cm,
  xlabel={délai d'obtention, en jours}, ylabel={part cumulée, en \%},
  xmin=0, xmax=90, ymin=0, ymax=105,
  grid=major, grid style={gray!25},
  every axis plot/.append style={thick},
  legend style={at={(0.5,-0.36)}, anchor=north, legend columns=2, font=\footnotesize,
    % SÉPARATION ENTRE LES DEUX COLONNES DE LA LÉGENDE. Sans elle, la clé de la seconde entrée
    % touchait le libellé de la première, et les deux se lisaient d'un seul tenant. Le défaut se
    % voyait sur les quatre graphiques à légende en deux colonnes, et sur eux seuls.
    /tikz/every even column/.append style={column sep=1.2em},
    draw=gray!40},
]
\addplot[color=black, mark=none] table[x=jours, y=artefact, col sep=comma]
  {\serie{delais-repartition}};
\addlegendentry{délai égal à la marge disponible}
\addplot[color=black, dashed, mark=none] table[x=jours, y=complement, col sep=comma]
  {\serie{delais-repartition}};
\addlegendentry{complément}
\end{axis}
\end{tikzpicture}
\caption{Répartition cumulée du délai d'obtention, sur les deux parts de la partition}\label{fig:delai-repartition}
\end{figure}

\conclusion{C-06}\textbf{Les deux courbes ont presque le même centre.} La médiane vaut
\chiffre{delai-mediane-ensemble} jours sur l'ensemble et \chiffre{delai-mediane-complement} sur le
complément. Le plafonnement déplace la queue de la distribution et laisse son centre en place. Les
médianes restent donc utilisables malgré lui.

Le délai se lit aussi par activité, sur les deux populations (cf. tableau~\ref{tab:delai-activite}).

\begin{table}[htbp]
\centering
\caption{Délai d'obtention par activité, sur les rendez-vous honorés à délai strictement positif}\label{tab:delai-activite}
\footnotesize
\pgfplotstabletypeset[
  col sep=comma,
  columns/code/.style={column name={Activité}, string type},
  columns/honores/.style={column name={Honorés}, int detect, 1000 sep={\,}},
  columns/mediane/.style={column name={Médiane}, fixed, fixed zerofill, precision=1},
  columns/decile9/.style={column name={9\textsuperscript{e} décile}, fixed, fixed zerofill,
    precision=1},
  columns/complement/.style={column name={Complément}, int detect, 1000 sep={\,}},
  columns/mediane_complement/.style={column name={Médiane}, fixed, fixed zerofill, precision=1},
  columns/decile9_complement/.style={column name={9\textsuperscript{e} décile}, fixed,
    fixed zerofill, precision=1},
  every head row/.style={before row=\toprule, after row=\midrule},
  every last row/.style={after row=\bottomrule},
]{\serie{delais-par-activite}}
\end{table}

Les délais y sont en jours. Les trois premières colonnes portent l'ensemble, les trois dernières le
complément. Les codes d'activité suivent la nomenclature nationale des spécialités médicales
\cite{s30} : \chiffre{code-activite-cardiologie} cardiologie, \chiffre{code-activite-dermatologie}
dermatologie, \chiffre{code-activite-gyneco-obstetrique} gynéco-obstétrique,
\chiffre{code-activite-medecine-generale} médecine générale, \chiffre{code-activite-medecine-interne}
médecine interne, \chiffre{code-activite-ophtalmologie} ophtalmologie,
\chiffre{code-activite-oto-rhino-laryngologie} oto-rhino-laryngologie et
\chiffre{code-activite-pediatrie} pédiatrie.

\conclusion{C-07}\textbf{Le classement des activités par délai reste le même d'une population à
l'autre.} L'écart entre les deux médianes d'une activité ne dépasse jamais un jour. La médiane la
plus haute vaut \chiffre{delai-rapport-extremes} fois la plus basse. Les paramètres fixent la
médiane de chaque activité, et le classement en découle.

\conclusion{C-08}\textbf{La part des rendez-vous pris pour le jour même dépasse son paramètre, et la
marge explique l'écart.} Elle vaut \chiffre{part-jour-meme-arrondi}~\% des honorés, pour un paramètre
de \chiffre{injecte-part-jour-meme}. Parmi eux, \chiffre{jour-meme-marge-nulle} rendez-vous, soit
\chiffre{jour-meme-part-marge-nulle-arrondi}~\%, ont une marge nulle ou négative. Leur fiche patient
date du jour du rendez-vous, et le générateur les place au jour même sans consulter le paramètre. Au
poste, un rendez-vous observé portait une seconde d'écart entre sa création et son
heure\releve{REL-RDV.C02} : le phénomène existe, et sa fréquence reste à mesurer.

\section{L'absentéisme et les annulations}

L'absence et l'annulation se distinguent nettement, et la relation entre délai et absence change
selon la façon de la mesurer. Le fait de rendez-vous porte \chiffre{source-rendez-vous-lignes} lignes
en quatre états\releve{REL-RDV.R11} : \chiffre{honores-total} honorés, \chiffre{absences-total}
absences, \chiffre{annulations-total} annulations et \chiffre{rdv-en-instance} encore ouverts. Ce
dernier état ne correspond à aucune des \chiffre{rdv-etats-signification} significations établies, et il ne reçoit aucun drapeau.

\conclusion{C-09}\textbf{Le taux d'absentéisme vaut \chiffre{taux-absenteisme-global-arrondi}~\% et le
taux d'annulation \chiffre{taux-annulation-global-arrondi}~\%.} Les deux se rapportent à l'ensemble
des rendez-vous, \chiffre{rdv-etats} états compris. Le taux d'annulation retrouve son paramètre,
\chiffre{injecte-taux-annulation}. Le classement des activités par absentéisme est fixé de la même
manière.

\conclusion{C-10}\textbf{L'annulation et l'absence se distinguent par deux champs.} Une annulation
porte un agent d'annulation\releve{REL-RDV.C07} et une date d'annulation\releve{REL-RDV.C08}.
Sur les \chiffre{annulations-total} annulations, \chiffre{annulations-avec-agent} portent les deux,
et \chiffre{absences-avec-agent-annulation} absence en porte un. Cette distinction vient du modèle de
données et non d'un paramètre : c'est l'une des rares conclusions que la chaîne établit vraiment.

Par activité, la part des absences se lit en regard du délai médian d'obtention (voir
figure~\ref{fig:delai-absenteisme}).

\begin{figure}[htbp]
\centering
\begin{tikzpicture}
\begin{axis}[
  width=11cm, height=4.6cm,
  xlabel={délai médian d'obtention, en jours}, ylabel={part des absences, en \%},
  xmin=0, xmax=52, ymin=0,
  grid=major, grid style={gray!25},
  legend style={at={(0.5,-0.36)}, anchor=north, legend columns=2, font=\footnotesize,
    % SÉPARATION ENTRE LES DEUX COLONNES DE LA LÉGENDE. Sans elle, la clé de la seconde entrée
    % touchait le libellé de la première, et les deux se lisaient d'un seul tenant. Le défaut se
    % voyait sur les quatre graphiques à légende en deux colonnes, et sur eux seuls.
    /tikz/every even column/.append style={column sep=1.2em},
    draw=gray!40},
]
\addplot[only marks, mark=*, mark size=2pt, color=black] table[x=mediane, y=absences,
  col sep=comma] {\serie{absenteisme-par-activite}};
\addlegendentry{ensemble}
\addplot[only marks, mark=square, mark size=2.4pt, color=black] table[x=mediane_complement,
  y=absences_complement, col sep=comma] {\serie{absenteisme-par-activite}};
\addlegendentry{complément}
\end{axis}
\end{tikzpicture}
\caption{Part des absences et délai médian d'obtention, par activité}\label{fig:delai-absenteisme}
\end{figure}

Chaque point y est une activité. Le complément retire des honorés et garde toutes les absences : son
niveau d'absences se compare mal à celui de l'ensemble, et seule la pente des deux nuages se compare.

\conclusion{C-11}\textbf{La relation entre délai et absentéisme change de signe selon la façon de la
mesurer.} Les trois mesures portent sur \chiffre{relation-effectif-ensemble} rendez-vous à issue
connue et à délai positif. Entre activités, sur \chiffre{activites-distinctes} points, la corrélation vaut
\chiffre{relation-inter-ensemble-arrondi} : les spécialités les plus attendues sont celles où les
patients se présentent le plus. À l'intérieur de chaque activité, après centrage, elle vaut
\chiffre{relation-intra-ensemble-arrondi} : à spécialité égale, un rendez-vous lointain est plus
souvent manqué. Sans distinguer les activités, elle vaut \chiffre{relation-brute-ensemble-arrondi}.

Cet écart est un effet de composition. La mesure brute mélange les deux autres et suit surtout la
mesure interne, parce que l'issue varie peu d'une activité à l'autre et beaucoup à l'intérieur de
chacune. Lue seule, la mesure entre activités ferait conclure qu'attendre ne fait pas manquer un
rendez-vous. Les deux relations viennent des paramètres : la corrélation entre leurs deux tables vaut
\chiffre{injecte-correlation-parametres-arrondi}, et la mesure la retrouve.

\conclusion{C-12}\textbf{Sur le complément, les trois mesures se rapprochent de la relation
injectée.} Sur \chiffre{relation-effectif-complement} rendez-vous, la mesure entre activités passe à
\chiffre{relation-inter-complement-arrondi}, la mesure interne à
\chiffre{relation-intra-complement-arrondi} et la mesure brute à
\chiffre{relation-brute-complement-arrondi}. Les délais égaux à la marge diluaient donc le lien sans
le porter.

\section{Les urgences}

Les urgences reçoivent un flux régulier, dont la gravité et l'orientation suivent la calibration. Le
service compte \chiffre{passages-urgences-produits} passages sur la période, soit
\chiffre{flux-urgences-moyenne} par jour.

\conclusion{C-13}\textbf{Le profil horaire des urgences forme un renflement de jour.} La pointe est à
\chiffre{urgences-heure-pointe} heures et le creux à \chiffre{urgences-heure-creux} heures, avec un
rapport de \chiffre{urgences-rapport-pointe-creux-arrondi} entre les deux (voir
figure~\ref{fig:passages-heure}). Le flux programmé, lui, a des heures sans aucun passage.

Les passages et le délai avant prise en charge médicale se ventilent par niveau de tri (cf.
tableau~\ref{tab:urgences-tri}).

\begin{table}[htbp]
\centering
\caption{Passages et délai avant prise en charge médicale, par niveau de tri}\label{tab:urgences-tri}
\footnotesize
\pgfplotstabletypeset[
  col sep=comma,
  columns/niveau/.style={column name={Niveau de tri}, int detect},
  columns/passages/.style={column name={Passages}, int detect, 1000 sep={\,}},
  columns/part/.style={column name={Part (\%)}, fixed, fixed zerofill, precision=2},
  columns/mediane/.style={column name={Médiane (min)}, fixed, fixed zerofill, precision=1},
  columns/decile9/.style={column name={9\textsuperscript{e} décile (min)}, fixed, fixed zerofill,
    precision=1},
  every head row/.style={before row=\toprule, after row=\midrule},
  every last row/.style={after row=\bottomrule},
]{\serie{urgences-par-niveau}}
\end{table}

Les \chiffre{niveaux-tri} codes de niveau s'y affichent sans libellé. La source qui décrit la nomenclature distingue
trois groupes de gravité \cite{s12}, et les données en portent \chiffre{niveaux-tri}.

\conclusion{C-14}\textbf{Le délai avant prise en charge croît avec le numéro du niveau de tri.} Les
cibles de délai, adossées à l'échelle de tri publiée \cite{s14}, valent
\chiffre{cible-delai-niveau-1}, \chiffre{cible-delai-niveau-2}, \chiffre{cible-delai-niveau-3},
\chiffre{cible-delai-niveau-4} et \chiffre{cible-delai-niveau-5} minutes. Chaque médiane reste sous
la cible de son niveau, et chaque neuvième décile la dépasse. Le taux de respect des cibles est
lui-même un paramètre : l'indicateur de conformité affiche un réglage et non une performance.

\conclusion{C-15}\textbf{La plupart des patients rentrent à domicile.} Le retour à domicile concerne
\chiffre{urgences-part-retour-domicile-arrondi}~\% des passages, le transfert
\chiffre{urgences-part-transfert-arrondi}~\% et l'hospitalisation
\chiffre{urgences-part-hospitalisation-arrondi}~\%. Le règlement intérieur nomme les cinq issues du
passage \cite{s27}. Le mode d'arrivée le plus fréquent porte
\chiffre{urgences-part-mode-principal-arrondi}~\% des passages, et son code reste sans libellé faute
de source.

\conclusion{C-16}\textbf{Les deux niveaux de tri les moins graves réunissent
\chiffre{urgences-part-consultation-ordinaire-arrondi}~\% des passages.} Ce sont les niveaux dont la
cible de délai est la plus longue, et ils relèvent d'une consultation ordinaire. Ce constat est établi
par les sources, et non par cette mesure. Le recueil national compte
\chiffre{source-part-non-urgentes}~\% de consultations non urgentes dans le public \cite{s12}. Une
étude dans un hôpital provincial mesure \chiffre{source-part-non-appropriees-arrondi}~\% de
consultations non appropriées \cite{s13}, au sens du recours et non du tri.

\section{La facturation et le recouvrement}

Le montant facturé suit la nature des épisodes plutôt que leur nombre, et le restant dû est ancien.
Le tableau~\ref{tab:montants-episode} le ventile par type d'épisode.

\conclusion{C-17}\textbf{Le montant facturé ne suit pas le nombre de factures.} Le type d'épisode le
moins nombreux porte la majorité du montant, et le plus nombreux un quart. Les
\chiffre{facturation-factures} factures portent \chiffre{facturation-montant-total} dirhams. La part
prise en charge par un organisme vaut \chiffre{facturation-part-organisme-arrondi}~\%, et la part du
patient \chiffre{facturation-part-patient-arrondi}~\%. Les deux parts sont mesurées séparément.

\conclusion{C-18}\textbf{Un taux de recouvrement se lit avec sa population.} Les
\chiffre{recouvrement-creances} créances nées portent \chiffre{recouvrement-montant-du} dirhams. Leur
taux de recouvrement vaut \chiffre{recouvrement-taux-global-arrondi}~\% et va de
\chiffre{recouvrement-taux-minimal-arrondi}~\% à \chiffre{recouvrement-taux-maximal-arrondi}~\% selon
le débiteur. Le taux d'encaissement est une autre grandeur : il rapporte les
\chiffre{encaissements-montant} dirhams encaissés au montant facturé, et vaut
\chiffre{taux-encaissement-arrondi}~\%.

L'ancienneté des créances non soldées pèse sur le restant dû (voir
figure~\ref{fig:anciennete-creances}).

\begin{figure}[htbp]
\centering
\begin{tikzpicture}
\begin{axis}[
  width=11cm, height=4.2cm,
  xlabel={ancienneté de la créance, en jours}, ylabel={part du restant dû, en \%},
  xmin=0, ymin=0, ymax=105,
  grid=major, grid style={gray!25},
  every axis plot/.append style={thick},
]
\addplot[color=black, mark=none] table[x=jours, y=part, col sep=comma]
  {\serie{anciennete-creances}};
\end{axis}
\end{tikzpicture}
\caption{Part du montant restant dû selon l'ancienneté des créances}\label{fig:anciennete-creances}
\end{figure}

La courbe donne, pour chaque ancienneté, la part du restant dû portée par des créances plus
anciennes. Son dénominateur est le montant restant des seules créances non soldées.

\conclusion{C-19}\textbf{Le restant dû est ancien.} Sur \chiffre{recouvrement-montant-restant}
dirhams à recouvrer, \chiffre{creances-part-au-dela-un-an-arrondi}~\% viennent de créances d'un an et
plus. Les \chiffre{recouvrement-relances} relances émises aboutissent, en tout ou en partie, dans
\chiffre{recouvrement-taux-aboutissement-arrondi}~\% des cas. Ce rendement se rapporte aux relances,
et non aux créances.

\section{La qualité du fichier patient}

Le défaut de saisie se concentre sur peu de colonnes, la quarantaine reste rare, et une même
personne porte parfois deux fiches. Le taux de renseignement des colonnes incomplètes se lit rang par
rang (voir figure~\ref{fig:completude-colonnes}).

\begin{figure}[htbp]
\centering
\begin{tikzpicture}
\begin{axis}[
  width=11cm, height=4.2cm,
  xlabel={colonnes incomplètes, rangées par complétude croissante},
  ylabel={renseignement, en \%},
  xmin=0, ymin=0, ymax=105,
  grid=major, grid style={gray!25},
  every axis plot/.append style={thick},
]
\addplot[color=black, mark=*, mark size=1pt] table[x=rang, y=completude, col sep=comma]
  {\serie{completude-champs}};
\end{axis}
\end{tikzpicture}
\caption{Taux de renseignement des colonnes incomplètes, par complétude croissante}\label{fig:completude-colonnes}
\end{figure}

Son abscisse est un rang de colonne, et non un identifiant.

\conclusion{C-20}\textbf{L'incomplétude se concentre sur quelques colonnes.} Sur
\chiffre{qualite-colonnes} colonnes des couches aval, \chiffre{qualite-colonnes-incompletes} sont
incomplètes et \chiffre{qualite-colonne-vide} est entièrement vide. La courbe monte lentement puis se
redresse : un petit nombre de colonnes porte presque tout le défaut.

\conclusion{C-21}\textbf{La quarantaine reste marginale.} Sur \chiffre{qualite-lignes-examinees}
lignes examinées, \chiffre{quarantaine-lignes-total-analyse} ont été rejetées à l'ingestion. Le taux
le plus élevé d'une table vaut \chiffre{qualite-taux-quarantaine-maximal-arrondi}~\%, et
\chiffre{qualite-tables-sans-quarantaine} tables sur \chiffre{qualite-tables} n'en rejettent aucune.
Le tableau~\ref{tab:qualite-par-table} donne le détail par table.

\conclusion{C-22}\textbf{Une même personne peut porter plusieurs identifiants.} Le fichier porte
\chiffre{identifiants-distincts} identifiants patient pour \chiffre{personnes-estimees} personnes
estimées. L'écart vient de \chiffre{grappes-multiples} grappes de plus d'une fiche, dont aucune n'en
réunit plus de \chiffre{grappes-taille-maximale}. Deux fiches désignent la même personne pour
\chiffre{doublons-part-mesuree-arrondi}~\% des personnes. Par collision exacte,
\chiffre{collision-groupes-nom} groupes partagent nom et date de naissance, et
\chiffre{collision-groupes-piece} une pièce d'identité. Une source établit le phénomène \cite{s20},
et son ampleur ici vient d'un paramètre.

\section*{Synthèse}

Les flux, les délais, l'absentéisme, les urgences, la facturation et la qualité du fichier se
calculent tous depuis les faits de la couche analytique. La plupart des résultats retrouvent un
paramètre du générateur, ce qui montre que la chaîne sait le retrouver. La relation entre délai et
absentéisme change de signe selon la façon de la mesurer, et la distinction entre absence et
annulation tient au modèle de données. Le chapitre~\ref{chap:tableau-de-bord} présente le tableau de
bord qui met ces grandeurs à la disposition du service.
